\section{Periodic notation, matrix conventions, and units}
\label{sec:notation}

\subsection{Direct and reciprocal lattices}

Using the column convention, define the primitive lattice matrix as
\begin{equation}
  \mathbf A=(\mathbf a_1\;\mathbf a_2\;\mathbf a_3)
\end{equation}
and the cell volume as
\begin{equation}
  \Omega=\det\mathbf A>0.
  \label{eq:cell-volume}
\end{equation}
The reciprocal-lattice matrix is
\begin{equation}
  \mathbf B=2\pi\mathbf A^{-\Herm},
  \qquad
  \mathbf a_i\cdot\mathbf b_j=2\pi\delta_{ij}.
  \label{eq:reciprocal-lattice}
\end{equation}
A lattice translation is $\vT=\mathbf A\mathbf n$ with
$\mathbf n\in\mathbb Z^3$, and a reciprocal vector is
$\vG=\mathbf B\mathbf m$.

Integrals over one cell are denoted by $\int_\Omega\dd\mathbf r$.  Constants
in a Fourier pair depend on the FFT normalization used by an implementation;
here we adopt
\begin{equation}
  \rho(\vG)=\int_\Omega e^{-i\vG\cdot\mathbf r}\rho(\mathbf r)\dd\mathbf r,
  \qquad
  \rho(\mathbf r)=\frac{1}{\Omega}\sum_{\vG}
  e^{i\vG\cdot\mathbf r}\rho(\vG).
  \label{eq:fourier-pair}
\end{equation}

\subsection{Bloch AOs and k-point weights}

From an atom-centered Gaussian $\chi_\mu$, define a Bloch AO by
\begin{equation}
  \phi_{\mu\vk}(\mathbf r)
  =\frac{1}{\sqrt{N_{\mathrm c}}}
   \sum_{\vT}e^{i\vk\cdot\vT}
   \chi_\mu(\mathbf r-\vT),
  \label{eq:bloch-ao}
\end{equation}
where $N_{\mathrm c}$ is the number of primitive cells in the Born--von Karman
supercell; it cancels from the final intensive quantities.  A Bloch orbital is
\begin{equation}
  \psi_{p\vk}(\mathbf r)
  =\sum_\mu C_{\mu p}^{\vk}\phi_{\mu\vk}(\mathbf r).
  \label{eq:bloch-mo}
\end{equation}

General weights are normalized as $\sum_\vk w_\vk=1$.  Regular meshes in this
package use equal weights, $w_\vk=1/N_k$.  Thus, an explicit factor $1/N_k$
below denotes the equal-weight convention.  A Gamma-only calculation is not a
separate scalar theory; it is the same tensor convention with $N_k=1$ and
$\vk=\mathbf0$.

The AO overlap and one-electron matrices are
\begin{align}
  S_{\mu\nu}^{\vk}
  &=\langle\phi_{\mu\vk}|\phi_{\nu\vk}\rangle,\\
  h_{\mu\nu}^{\vk}
  &=\langle\phi_{\mu\vk}|\hat h|\phi_{\nu\vk}\rangle,
  \label{eq:ao-matrices}
\end{align}
and the MO coefficients satisfy
\begin{equation}
  C^{\vk\dagger}S^{\vk}C^{\vk}=I.
  \label{eq:mo-orthonormality}
\end{equation}

\subsection{Restricted density matrix}

For closed-shell occupations $f_{p\vk}\in\{0,2\}$, the AO density is
\begin{equation}
  D_{\mu\nu}^{\vk}
  =\sum_p f_{p\vk}C_{\mu p}^{\vk}C_{\nu p}^{\vk *}
  =2\sum_i C_{\mu i}^{\vk}C_{\nu i}^{\vk *}.
  \label{eq:restricted-density}
\end{equation}
This definition already contains the occupation factor of two.  Whether an
additional factor of two is required in an energy or response contraction
therefore depends on whether an occupied-only tensor is being used in place of
$D$.

The matrix-trace convention is
\begin{equation}
  \Tr[D^{\vk}X^{\vk}]
  =\sum_{\mu\nu}D_{\nu\mu}^{\vk}X_{\mu\nu}^{\vk}.
  \label{eq:trace-convention}
\end{equation}
At complex k-points, $D^{\vk}$ and a Hermitian operator are Hermitian, but an
individual derivative block need not be.  The real part is taken only for the
final scalar; intermediate imaginary components must not be discarded
arbitrarily.

\subsection{Indices and derivative notation}

\begin{center}
\begin{tabular}{cl}
\toprule
Index & Meaning \\
\midrule
$\mu,\nu,\lambda,\sigma$ & primary AO \\
$P,Q$ & auxiliary AO \\
$i,j$ & occupied MO \\
$a,b$ & virtual MO \\
$p,q$ & arbitrary MO \\
$A,B$ & atom \\
$x,y$ & Cartesian component \\
$\vk,\vk'$ & sampled k-point \\
\bottomrule
\end{tabular}
\end{center}

A total derivative with respect to $R_{Ax}$ is written as a superscript $Ax$
or as $\partial_{Ax}$; for example,
$h^{Ax}=\partial h/\partial R_{Ax}$.  Because the AO basis moves with the
nuclei, a matrix derivative must be distinguished from the orbital-coefficient
response.  The cell is fixed, so $\mathbf A^{Ax}=0$, and the reciprocal vectors
and discrete mesh are also fixed.

\subsection{Canonical units}

\begin{center}
\begin{tabular}{lll}
\toprule
Quantity & Canonical unit & Comment \\
\midrule
Energy & Hartree (Ha) & public energy \\
Position/lattice & Bohr & PySCF internal unit \\
Nuclear gradient & Ha/Bohr & $\dd E/\dd R$ \\
Force & Ha/Bohr & $-\dd E/\dd R$ \\
k-point & Bohr$^{-1}$ & Cartesian reciprocal coordinate \\
Plane-wave cutoff & Ha & $G_{\max}^2/2$ \\
\bottomrule
\end{tabular}
\end{center}

Even when an input XYZ file is given in angstrom, the coordinates are in Bohr
after the cell is built.  Converting a gradient array to eV/angstrom requires
both energy and length conversions.  Changing the sign is not a unit
conversion; it converts a gradient into a force.
